\documentclass[11pt]{article}
\usepackage{amsmath,amssymb,amsthm}
\newtheorem{theorem}{Theorem}
\newtheorem{lemma}{Lemma}
\title{Compact Joint Zero Families from Chebyshev Fibers}
\author{}
\date{November 24, 2026}
\begin{document}
\maketitle

\paragraph{Status and scope.}
This is a complete written construction of many exact zero representations
on a one-hot grid, together with their full first-order perturbation
space. All denominators satisfy strict reverse admissibility on the
whole Boolean cube. It is not a local compiler: no theorem below says
that the union of the perturbation spaces contains every sign table.
Consequently it establishes neither accepted asymptotic target. It has
not been formalized in Lean. The companion script checks rational
quadratic and cubic instances and their cube budgets exactly.

Let $T_k,U_{k-1}$ denote the Chebyshev polynomials, normalized by
$T_k(\cos\theta)=\cos(k\theta)$ and
$T_k'=kU_{k-1}$. For $-1<t<1$, let
$\xi_1(t)<\cdots<\xi_k(t)$ be the distinct roots of $T_k(X)-t$.
All these roots lie in $(-1,1)$.

\begin{theorem}[A joint zero family]
Fix positive integers $r,s,k$, values $t_j\in(-1,1)$, and a sufficiently
small $\varepsilon>0$. There is an open interval $I$ containing $1$
such that, for any $x_1,\ldots,x_r\in I$ and any independent
permutations $\pi_j\in S_k$, the grid identity
\[
 \sum_{h=1}^k\frac{1}{x_i-\varepsilon\xi_{\pi_j(h)}(t_j)}
 =\frac{(k/\varepsilon)U_{k-1}(x_i/\varepsilon)}
 {T_k(x_i/\varepsilon)-t_j}
 \tag{1}
\]
holds. Every denominator displayed in (1) is the restriction of an
affine denominator with strictly negative slopes and positive values
on the entire Boolean $(r+s)$-cube. The numerator on the right is
also the restriction of an affine function. Thus $k+1$ strict reverse
heads give a grid-zero representation for every permutation tuple.

If $s\geq1$, fixing $\pi_1$ leaves $(k!)^{s-1}$ distinct denominator
multisets on the grid. Distinctness here is modulo a simultaneous
relabeling of the first $k$ heads, with the baseline head distinguished.
\end{theorem}

\begin{proof}
Logarithmically differentiating the factorization of $T_k(X)-t$ gives
\[
 \sum_{h=1}^k\frac1{X-\xi_h(t)}
 =\frac{T_k'(X)}{T_k(X)-t}.
\]
Substitution $X=x_i/\varepsilon$ proves (1).

Here are direct full-cube extensions, including strict slope padding.
First set $x_i=1$ for all $i$. For each $h$, write
$v_{hj}=\varepsilon\xi_{\pi_j(h)}(t_j)$. The denominator
$1-v_{hj}$ on the grid is extended by
\[
 B_h(z,w)=1+\varepsilon+2\eta
       -\eta\sum_i z_i-\sum_j(\varepsilon+v_{hj}+\eta)w_j.
 \tag{2}
\]
Every slope is negative. Since $|v_{hj}|<\varepsilon$, its value at
the all-one vertex is greater than
$1-(2s-1)\varepsilon-(r+s-2)\eta$.
Choose, for example, $0<\varepsilon<1/(8s)$ and then $\eta>0$
small enough that this is positive. The all-one vertex is the minimum.

Put $C=T_k(1/\varepsilon)$ and $t_* =\min_j t_j$.
The baseline denominator has the extension
\[
 B_0(z,w)=C-t_*+2\eta_0
 -\eta_0\sum_i z_i-\sum_j(t_j-t_*+\eta_0)w_j.
 \tag{3}
\]
It has the required grid value $C-t_j$ and strict negative slopes.
Because $T_k(u)\geq u$ for $u\geq1$, our choice gives $C>8s$.
Before the padding its all-one value is
$C-t_*-\sum_j(t_j-t_*)>C-1-2s>0$.
Choose $\eta_0>0$ sufficiently small.

Now let $x_i$ vary near $1$. In (2), replace the coefficient $\eta$
of $z_i$ by $\eta+1-x_i$. In (3), replace $\eta_0$ by
$\eta_0+C-T_k(x_i/\varepsilon)$. These changes give exactly the
desired grid denominators. At $x_i=1$ all slope inequalities and
all cube positivity inequalities are strict. Continuity therefore
preserves them for $x_i$ in a sufficiently small common interval $I$.
There are finitely many permutation tuples, so $I$ may be chosen
uniformly for all of them. In fact (2) already gives uniform slack
independent of the tuple. Finally
$\sum_i (k/\varepsilon)U_{k-1}(x_i/\varepsilon)z_i$ is an affine
numerator with the stated grid values.

For distinctness, the first column has distinct pole labels. A
simultaneous head relabeling preserving this column is therefore the
identity once $\pi_1$ is fixed. Every later column then determines its
permutation uniquely from its denominator values.
\end{proof}

\paragraph{The precise perturbation family.}
Set
\[
 D_{hij}=x_i-\varepsilon\xi_{\pi_j(h)}(t_j),\quad
 D_{0ij}=T_k(x_i/\varepsilon)-t_j,\quad
 N_i=(k/\varepsilon)U_{k-1}(x_i/\varepsilon).
\]
Take the first $k$ base numerators to be $1$, and the signed baseline
numerator to be $-N_i$. Independent affine numerator and denominator
perturbations on the grid have values $a_{hi}+b_{hj}$ and
$p_{hi}+q_{hj}$, respectively. Their first-order grid score is
\[
 L_{ij}=\sum_{h=1}^k\left(
 \frac{a_{hi}+b_{hj}}{D_{hij}}
 -\frac{p_{hi}+q_{hj}}{D_{hij}^2}\right)
 -\frac{a_{0i}+b_{0j}}{D_{0ij}}
 +\frac{N_i(p_{0i}+q_{0j})}{D_{0ij}^2}.
 \tag{4}
\]
All the parameters in (4) are free real numbers. To verify this
without a hidden affine restriction, realize each row-plus-column
table as $\sum_i a_{hi}z_i+\sum_j b_{hj}w_j$; do the same for the
denominator direction. The base point is interior to the admissible
denominator cone, so sufficiently small perturbations of either sign
remain admissible.

\begin{lemma}[Strict tangent signs transfer]
If, for some tuple of permutations and some values of the parameters
in (4), one has $\sigma(i,j)L_{ij}>0$ at every grid point, then a
sufficiently small positive perturbation gives $k+1$ strict reverse
heads with those grid signs.
\end{lemma}
\begin{proof}
The base score vanishes by (1). The perturbed score is
$\delta L_{ij}+O(\delta^2)$ at each grid point. There are finitely
many grid points, and their first-order margins are all strict, so
one common positive $\delta$ works. The preceding interior argument
preserves whole-cube admissibility.
\end{proof}

\paragraph{What this does and does not evade.}
With only row numerator perturbations, (4) is a Cauchy assignment
readout. Its fixed-residue butterfly subfamily is covered by the
first-layer dominance obstruction in
\texttt{butterfly-rectangle-decoder.tex}. Independent column numerator
perturbations let the residues depend on the entire target column and
its selected permutation. Independent denominator perturbations also
add squared Cauchy kernels. Those choices are outside the fixed-residue
hypothesis, but escaping that hypothesis is not a coverage proof.

The zero identity is only on the one-hot grid. Axis-zero conditions
must still be imposed using an audited repair with its stated head
cost, or by a new identity. Neither the permutation count nor the
dimension of (4) proves that every sign table is covered. The missing
statement remains a quantitative joint sign-coverage theorem with
$k=o(r+s)$.
\end{document}
